
\section{$s_1,s_2$ 未知时是否还能只用相对位置信息处理}

\subsection{原文对 $S$ 的假设}

原文第 2 节明确写道：
\begin{quote}
“Although the system matrix $S$ is known, each vehicle lacks direct access
to the target's position $x_0$ and velocity $v_0$.”
\end{quote}
即原文假设\textbf{模型矩阵 $S$ 是已知的}，未知的是目标状态 $x_0,v_0$（由相对位置测量与内部模型推断）。

更关键的是，$S$ 直接出现在补偿器 (11) 的矩阵 $G_1$ 中：
\[
G_1 = \begin{bmatrix}
-l_1 & 1 & 0 & 0\\
-k_1-l_2 & -k_2 & -k_3 & -k_4\\
0 & 0 & 0 & 1\\
0 & 0 & -s_1 & -s_2
\end{bmatrix},\qquad
G_2 = \begin{bmatrix} l_1 \\ l_2 \\ 0 \\ 1\end{bmatrix}.
\]
其下右 $2\times2$ 块正是
\[
\begin{bmatrix}0&1\\ -s_1&-s_2\end{bmatrix}=S.
\]
这正是内部模型原理的要求：控制器中必须“包含”外系统动力学 $S$。

\subsection{为何 $s_1,s_2$ 未知时原算法直接失效}

内部模型原理要求控制器嵌入 $S$；原文的 $G_1$ 把 $S$ 显式写进了下右块。
若 $s_1,s_2$ 未知，则 $G_1$ 无法构造，固定结构的补偿器 (11) 根本没有定义——
这并不是“测量不足”的问题，而是“模型未知导致内部模型无法搭建”。
因此，\textbf{原文的固定结构算法不能处理未知的 $s_1,s_2$}。

\subsection{能否改用“仅相对位置”来处理未知 $S$？——需要自适应内部模型}

未知外系统参数时，标准输出调节框架需升级为
\textbf{自适应输出调节}（adaptive output regulation，
Serrani--Isidori--Marconi 2001；Chen \& Huang 2015）：
把控制器中那一份 $S$ 改为\textbf{在线估计}的内部模型副本，即
\[
\hat S(t)=\begin{bmatrix}0&1\\ -\hat s_1(t)&-\hat s_2(t)\end{bmatrix},
\]
其中 $\hat s_1,\hat s_2$ 由自适应律根据调节误差（相对位置误差 $e_i=x_i-x_0$）
实时更新。这样 $S$ 不再需要预先已知。

关于“仅相对位置”：
\begin{itemize}
  \item 原补偿器本来只接收位置型相对测量 $e_i=x_i-x_0$（不含速度），
    目标状态 $v_0$ 也是由内部模型推断的；
  \item 相对位置误差里含有外系统信息（经由 $x_0$ 进入），
    因此只要“测量—内部模型”对可检测（detectable），
    且外系统被充分激励，\textbf{仅相对位置就足以驱动自适应内部模型}。
\end{itemize}

\subsection{可行性条件}

自适应方案要成立，一般需满足：
\begin{enumerate}
  \item 被控对象 $(A,B)$ 可稳、$(C,A)$ 可检测——
    本例为二阶积分器、测量位置，显然满足。
  \item 外系统的\textbf{结构}已知（二阶），仅参数未知——
    这样才能确定自适应内部模型的结构与维数。
    若连“二阶”这一结构都未知，则须先辨识结构，难度更大。
  \item \textbf{持续激励（persistent excitation）}：
    自适应律要收敛到真值，需调节误差对外系统频率充分激励。
    对原文“时变速度目标”，若 $S$ 的特征值在虚轴或右半平面
    （振荡或增长型速度），目标自然提供持续激励，参数可收敛；
    若目标为常速（$S=0$），则无激励，无法辨识 $S$（但此时 $S$ 本就是已知的零矩阵）。
  \item 自适应后的对 $(G_2C, G_1(\hat S))$ 可检测。
\end{enumerate}

\subsection{与原文场景的对应}

原文场景是 “target with time-varying velocity”，
即 $S$ 通常具虚轴或正实部特征值——这正是“自然持续激励”的情形。
因此，\textbf{改用自适应内部模型后，仅凭相对位置测量即可在未知 $s_1,s_2$ 下实现合围与速度同步}；
而原文“固定 $G_1$”算法本身不行。

\subsection{另一条路线：目标运动估计器（文献 [22,23]）}

原文引言提到 [22,23] 用有限时间估计器观测目标运动。
若 $S$ 未知，也可先由相对测量用估计器重建目标状态（及/或 $S$），
再把估计出的 $S$ 当作“已知”送入围捕控制器。
这等价于把“未知 $S$”交给估计器处理，控制器仍沿用已知（估计值）$S$ 的内部模型。
这是另一类算法，而非原文的固定内部模型算法。

\subsection{结论}

\begin{itemize}
  \item 原文的“固定结构”算法\textbf{不能}处理未知 $s_1,s_2$：
    $S$ 直接写进 $G_1$，且原文明确假设 $S$ 已知。
  \item 在“仅相对位置 + $S$ 未知”下，需要把算法改为
    \textbf{自适应内部模型}（或“目标运动估计器 + 围捕控制器”）方案。
    相对位置测量本身足以驱动自适应/估计，但需满足
    可检测性与持续激励条件；
    对原文的“时变速度目标”场景，持续激励自然满足，故可行。
  \item 简言之：未知 $s_1,s_2$ 改变的是“模型是否已知”，
    而非“测量是否充分”——所以补救办法是引入自适应/估计，
    而不是增加测量通道。
\end{itemize}
